\section{Global signs}\label{sec:signs}

We now prove the global signs for the exact functions supplied by
Proposition~\ref{interpolation:solution}.  The strict bounds will also
identify their zeros within the respective sign domains.  Write
\[
  \sigma_1=-1,\qquad \sigma_2=1,\qquad
  \delta=3\cdot10^{-5}.
\]
The difference between each exact coefficient list and its tabulated list has
norm less than $\delta$.  Their fixed coefficients at zero and their constants
$C_i$ agree.  Tildes will always denote the functions constructed from the
finite tabulated lists, so that
$\widetilde H_i=\widetilde p_i+K_{\widetilde p_{3-i}}$.
We use the spectral envelopes of Section~\ref{sec:interpolation} throughout.

\subsection{Taylor quotients on half-gaps}

Set $\mathcal N_0=\{0\}\cup\mathcal N$.  For each
$m\in\mathcal N_0\cap[0,40]$, let $m^+$ be its successor in
$\mathcal N_0$ and include
\[
  y=\frac{m^+-m}{2}.
\]
If $m>0$, also include $y=(m^--m)/2$, where $m^-$ is its predecessor.
Thus $y$ is signed, and the associated half-gap is
\[
  I_{m,y}=\{m+yv:0\le v\le1\}.
\]
For function index $2$ we use every such interval, with order
$\nu=0$ at $m=0$ and $\nu=2$ otherwise.  For function index $1$ we use only
those intervals contained in $[1,\infty)$, with $\nu=1$ at $m=1$ and
$\nu=2$ otherwise.  In particular, the only odd order occurs on the
rightward half-gap from $1$.
There are $22$ centers from $0$ through $40$.  They give $43$ half-gaps for
index $2$ and, after removing the two intervals not contained in $[1,\infty)$, $41$ for index
$1$, for a total of $84$ certified intervals.

Fix one of the specified triples $(i,m,y)$.  For $s=m+u$ in its interval,
we will prove positivity of
\[
 \frac{\sigma_i H_i(m+u)}{u^\nu}.
\]
Here $u^0=1$, and for $\nu>0$ the quotient at $u=0$ is interpreted by
continuity.  The interpolation conditions make the derivatives of $H_i$
of orders below $\nu$ vanish at $m$, so this extension exists.  Positivity
of the quotient gives the required sign: every order is even except for
the rightward half-gap from $1$, where $u\ge0$.
Figure~\ref{fig:half-gaps} shows the node-to-midpoint intervals on which
these quotients are studied.
\begin{figure}[htbp]
\centering
\setlength{\unitlength}{0.67cm}
\begin{picture}(14,2.35)
\put(0,1.2){\vector(1,0){13.4}}
\put(13.6,1.05){$s$}
\put(0,1.2){\circle*{.13}}
\put(1,1.2){\circle*{.13}}
\put(3,1.2){\circle*{.13}}
\put(4,1.2){\circle*{.13}}
\put(7,1.2){\circle*{.13}}
\put(9,1.2){\circle*{.13}}
\put(12,1.2){\circle*{.13}}
\put(13,1.2){\circle*{.13}}
\put(0,.65){\makebox(0,0){$0$}}
\put(1,.65){\makebox(0,0){$1$}}
\put(3,.65){\makebox(0,0){$3$}}
\put(4,.65){\makebox(0,0){$4$}}
\put(7,.65){\makebox(0,0){$7$}}
\put(9,.65){\makebox(0,0){$9$}}
\put(12,.65){\makebox(0,0){$12$}}
\put(13,.65){\makebox(0,0){$13$}}
\color{blue!55!black}
\put(.5,1.05){\line(0,1){.3}}
\put(2,1.05){\line(0,1){.3}}
\put(3.5,1.05){\line(0,1){.3}}
\put(5.5,1.05){\line(0,1){.3}}
\put(8,1.05){\line(0,1){.3}}
\put(10.5,1.05){\line(0,1){.3}}
\put(12.5,1.05){\line(0,1){.3}}
\put(4,1.8){\line(1,0){1.5}}
\put(4,1.65){\line(0,1){.3}}
\put(5.5,1.65){\line(0,1){.3}}
\put(6,1.8){\makebox(0,0)[l]{\small one half-gap}}
\end{picture}
\caption{Interpolation nodes (dots) and adjacent gap midpoints (short
marks) in the squared-radius coordinate. Each node-to-midpoint segment
is a half-gap. The finite sign certificates control a Taylor quotient
on each whole half-gap; a separate argument covers the remaining unbounded region.}
\label{fig:half-gaps}
\end{figure}


The tabulated function $\widetilde H_i$ approximates $H_i$ but need not
reproduce its exact zero and slope at a node.  We therefore compare the
exact quotient with the Taylor series of $\widetilde H_i$ after deleting
the terms of orders below $\nu$.  Put $\Delta_i=H_i-\widetilde H_i$.
Since $\widetilde H_i$ is entire and the lower derivatives of $H_i$ vanish,
this Taylor series gives
\begin{align}\label{signs:deleted-jets}
 &\frac{\sigma_i H_i(m+u)}{u^\nu}
 -\sigma_i\sum_{r=0}^{\infty}
       \frac{\widetilde H_i^{(r+\nu)}(m)}{(r+\nu)!}u^r
 \nonumber\\
 &\hspace{1cm}=
 \frac{\sigma_i}{u^\nu}
 \left(\Delta_i(m+u)-
       \sum_{j=0}^{\nu-1}\frac{\Delta_i^{(j)}(m)}{j!}u^j\right).
\end{align}
For $\nu=0$, the sum on the second line is empty and the right-hand side is
simply $\sigma_i\Delta_i(m+u)$.  For $\nu>0$, Taylor's integral remainder
gives the exact identity
\[
 \frac{\Delta_i(m+u)-\sum_{j=0}^{\nu-1}
             \Delta_i^{(j)}(m)u^j/j!}{u^\nu}
 =\frac{1}{(\nu-1)!}\int_0^1(1-t)^{\nu-1}
                \Delta_i^{(\nu)}(m+tu)\,\dd t.
\]
It holds for either sign of $u$ and extends continuously to $u=0$.
Since $m+tu$ stays in the same half-gap, the absolute value of the
right-hand side of \eqref{signs:deleted-jets} is at most
\begin{equation}\label{signs:coefficient-error}
  \frac{1}{\nu!}\sup_{s\in I_{m,y}}
       \abs{\Delta_i^{(\nu)}(s)}.
\end{equation}
The same bound applies at $u=0$ by continuity and for either orientation
of the half-gap.  Deleting the lower Taylor terms is essential here:
dividing the raw table error by $u^\nu$ could be singular at the node.

After writing $u=yv$, the first $29$ terms of the tabulated series in
\eqref{signs:deleted-jets} form a polynomial of degree at most $28$ in
$v$.  We bound it on $[0,1]$ using the standard range enclosure by Bernstein
coefficients \cite[Section~3]{TitiGarloff2019}, retaining its elementary
proof below.

\begin{lemma}[Finite Bernstein certificate]\label{signs:bernstein}
For every specified triple $(i,m,y)$, put
\begin{equation}\label{signs:power-coefficients}
  T_r=\sigma_i y^r
       \frac{\widetilde H_i^{(r+\nu)}(m)}{(r+\nu)!},
  \qquad 0\le r\le28.
\end{equation}
Then, for $0\le k\le28$,
\begin{equation}\label{signs:bernstein-bound}
  B_k:=\sum_{r=0}^k
       \frac{\binom{k}{r}}{\binom{28}{r}}T_r
  >\begin{cases}
      0.74,&m=0,\\
      0.18,&m=1,\\
      0.02,&m=3,4,\\
      0.009,&4<m\le40.
    \end{cases}
\end{equation}
\end{lemma}

\begin{proof}
The primary interval calculation in Appendix~\ref{sec:certificates} encloses
the exact derivatives of the tabulated functions in
\eqref{signs:power-coefficients}, using the quadrature assertion of
Lemma~\ref{cert:moment-error} for the analytic moment remainder.
Its exact grouped bounds~\eqref{cert:bernstein-table} imply the stated
thresholds.
\end{proof}

The relevance of this finite statement is the identity
\begin{equation}\label{signs:bernstein-identity}
  \sum_{r=0}^{28}T_r v^r
  =\sum_{k=0}^{28} B_k\binom{28}{k}v^k(1-v)^{28-k}.
\end{equation}
Indeed, for $0\le r\le k\le d=28$, the binomial identity
\[
 \frac{\binom{k}{r}}{\binom{d}{r}}\binom{d}{k}
 =\binom{d-r}{k-r}
\]
gives
\[
 \sum_{k=r}^{d}\frac{\binom{k}{r}}{\binom{d}{r}}
       \binom{d}{k}v^k(1-v)^{d-k}
 =v^r\sum_{l=0}^{d-r}\binom{d-r}{l}v^l(1-v)^{d-r-l}=v^r.
\]
Multiply by $T_r$, sum over $r$, and interchange the two finite sums to
obtain \eqref{signs:bernstein-identity}.
For $0\le v\le1$, the basis functions on the right are nonnegative and
sum to one.  Hence the lower bounds in
\eqref{signs:bernstein-bound} hold on the whole interval for the polynomial
on the left.  To obtain a lower bound for the exact quotient, it remains
to control the coefficient error in \eqref{signs:coefficient-error} and the
terms beyond degree $28$ in the tabulated series.

\subsection{Bounds for the two errors}

The density bounds and \eqref{interpolation:moment-bound}, together with
the separate coefficient error bounds $\delta$ for the two functions,
give, for $s\ge s_*\ge0$,
\begin{equation}\label{signs:perturbation-envelope}
 \abs{\Delta_i^{(\nu)}(s)}
 \le \delta\sum_{j=1}^6 M_j
       \left(\frac{(\pi t_j)^\nu}{6}
            +L_j(s_*)(\pi a_{j-1}^*)^\nu\right).
\end{equation}
There are no atom contributions because the constants $C_i$ have not
changed.  Substitution of the envelopes proves the following uniform
bounds for \eqref{signs:coefficient-error}:
\begin{equation}\label{signs:perturbation-caps}
 \frac{1}{\nu!}\sup_{I_{m,y}}\abs{\Delta_i^{(\nu)}}
 <\begin{cases}
    0.004,&m=0,\\
    0.011,&m=1,\\
    0.002,&m\ge3.
   \end{cases}
\end{equation}
Here one uses $s_*=0$ for $m=0$, $s_*=1/2$ for $m=1$ and both possible
orders, and $s_*=2$ for $m\ge3$.  For example the first bound is
$\delta\sum_{j=1}^6(M_j/6+M_jL_j(0))<0.004$.
All these are bounds on intervals, obtained from the derivative envelopes.

It remains to truncate the infinite series in
\eqref{signs:deleted-jets}.  Put $Y=\abs{y}$ and $q=29+\nu$.
The total variation of the direct spectral measure of either tabulated
function is less than $80$, since
\[
  2.29\sum_{j=1}^6\frac{M_j}{6}
       +0.017\sum_{j=0}^6 m_j^*<80.
\]
For the transformed part, use coefficient norm less than $2.3$ and
$\abs{C_i}\le0.017$.  Expanding each spectral exponential at $m$ gives
the following bound for the omitted quotient terms:
\begin{align}\label{signs:taylor-bound}
 \mathcal E(m,Y,\nu)
 := {}&\frac{Y^{-\nu}}{q!}
       \frac{1}{1-\pi a_0^*Y/(q+1)}
 \Bigg(80(\pi Y)^q
 \nonumber\\[-2pt]
 &\qquad+2.3\sum_{j=1}^6
       M_jL_j(m)(\pi a_{j-1}^*Y)^q
       +0.017\sum_{j=0}^6
       L_j^P(m)(\pi a_j^*Y)^q\Bigg).
\end{align}
Indeed the first omitted exponential term has degree $q$ and subsequent
absolute-value ratios are at most $\pi a_0^*Y/(q+1)<1$.
After division by $u^\nu$, all powers still have degree at least $29$,
so their maximum for $\abs{u}\le Y$ is bounded by evaluating these
absolute-value majorants at $Y$.  This proves
\eqref{signs:taylor-bound} also when $y<0$.

The five representative substitutions are
\begin{equation}\label{signs:taylor-representatives}
 (m,Y,\nu)=(0,\tfrac12,0),\ (1,1,1),\ (1,1,2),\ (3,1,2),\
           (4,\tfrac32,2).
\end{equation}
In each case \eqref{signs:taylor-bound} is less than $0.001$.
Increasing $m$ or decreasing $Y$ only decreases the bound.  The first
three cases handle centers $0$ and $1$, the fourth handles $m=3$, and the
last handles every center $m\ge4$.  Thus the truncation error is less
than $0.001$ on every required half-gap.

Combining \eqref{signs:bernstein-identity},
\eqref{signs:deleted-jets}, \eqref{signs:perturbation-caps}, and
\eqref{signs:taylor-bound}, we obtain
\begin{equation}\label{signs:finite-quotient}
  \frac{\sigma_i H_i(s)}{(s-m)^\nu}
  >\begin{cases}
     0.735,&m=0,\\
     0.168,&m=1,\\
     0.017,&m=3,4,\\
     0.006,&4<m\le40,
    \end{cases}
  \qquad s\in I_{m,y},
\end{equation}
with endpoint values interpreted by continuity.  Since the only odd
order has $s-m\ge0$, these inequalities imply the desired signs on all
of these intervals.  Their union is $[0,41.5]$ for index $2$ and
$[1,41.5]$ for index $1$: adjacent half-gaps meet at their common
midpoint, and the last interval runs from $40$ halfway to $43$.

\subsection{The unbounded region}

For $s\ge s_0:=41.5$, split the exact direct part as
\begin{equation}\label{signs:tail-split}
 \begin{aligned}
 p_i(s)&=P(s)R_{i,0}(s)+p_{i,1}(s),\\
 R_{i,0}(s)&=C_i+
   \sum_{\substack{n\in\mathcal N_0\\n\le40}}
   \left(\frac{c_{i,n}}{(s-n)^2}+\frac{d_{i,n}}{s-n}\right),\\
 \widetilde R_{i,0}(s)&=C_i+
   \sum_{\substack{n\in\mathcal N_0\\n\le40}}
   \left(\frac{\widetilde c_{i,n}}{(s-n)^2}
         +\frac{\widetilde d_{i,n}}{s-n}\right).
 \end{aligned}
\end{equation}
Here $p_{i,1}$ is the cardinal contribution from exact coefficients at
nodes greater than $40$.  The function $R_{i,0}$ uses the exact coefficients,
whereas
$\widetilde R_{i,0}$ uses the finite table; their constants and node-zero
coefficients agree.  All denominators in these two rational functions are
at least $1.5$.  Write the corresponding decomposition of the exact
function as
\[
 E_i:=p_{i,1}+K_{p_{3-i}},\qquad H_i=P R_{i,0}+E_i.
\]
For each $s\ge s_0$, the final comparison will use a nearest node $m\ge43$.
At such a node, the exact interpolation conditions and the double zero of
$P$ make $E_i$ have zero value and slope; we verify this at the comparison
step below.  The useful scale for both terms is therefore $(s-m)^2$.  We will bound
$\sigma_iR_{i,0}$ below, bound $P(s)$ below by a multiple of $(s-m)^2$,
and bound $|E_i(s)|$ above at the same scale through a uniform
second-derivative estimate.

We first show
\begin{equation}\label{signs:rational-tail}
  \sigma_i R_{i,0}(s)>0.011\qquad(s\ge s_0,\ i=1,2).
\end{equation}
For the tabulated coefficients put
$u_n^*=\sigma_i\widetilde c_{i,n}$,
$v_n^*=\sigma_i\widetilde d_{i,n}$, and $[w]_-:=\min\{w,0\}$.
Using
\[
 \frac{1}{s-n}=\frac1s+\frac{n}{s(s-n)}
\]
and discarding positive terms gives the following lower bound for
$\sigma_i\widetilde R_{i,0}(s)$, uniformly for $s\ge s_0$:
\begin{equation}\label{signs:rational-lower-sum}
 \sigma_i C_i+
 \frac{[\sum_{n\le40}v_n^*]_-}{s_0}
 +\sum_{n\le40}
    \left(\frac{[u_n^*]_-}{(s_0-n)^2}
          +\frac{[v_n^*]_-n}{s_0(s_0-n)}\right),
\end{equation}
where both sums include $n=0$.  Each negative reciprocal term is smallest
at $s=s_0$.  The exact rational sums in
\eqref{signs:rational-lower-sum} exceed $0.012$ for index $1$ and $0.016$
for index $2$.  To pass to the exact rational part, use $s-n\ge1.5>1$ and
the coefficient enclosure to obtain
\[
 \begin{split}
 \abs{R_{i,0}(s)-\widetilde R_{i,0}(s)}
 &\le\sum_{n\le40}\left(
   \frac{\abs{c_{i,n}-\widetilde c_{i,n}}}{(s-n)^2}
  +\frac{\abs{d_{i,n}-\widetilde d_{i,n}}}{s-n}\right)\\
 &<\frac{\delta}{s_0-40}=0.00002.
 \end{split}
\]
The finite lower sums therefore imply the asserted exact bound
\eqref{signs:rational-tail}.

The coefficient norm of the tabulated part with $n>40$ is less than
$0.00002$ for either function.  Hence the exact coefficient norm for
$p_{i,1}$ is less than $0.00005$.  The derivative envelopes imply
\begin{align}\label{signs:tail-second-derivative}
 \sup_{s\ge s_0}\abs{E_i''(s)}
 \le {}&0.00005\cdot77
       +2.3\sum_{j=1}^6 M_jL_j(s_0)(\pi a_{j-1}^*)^2
 \nonumber\\
       &+0.017\sum_{j=0}^6 L_j^P(s_0)(\pi a_j^*)^2
       <0.006.
\end{align}
For the direct part we used
$\sum_{j=1}^6 M_j(\pi t_j)^2/6<77$; for the transformed part we used the
all-node coefficient norm less than $2.3$ from
Proposition~\ref{interpolation:solution}.

It remains to give the quadratic lower bound for $P$.
\begin{lemma}\label{signs:P-lower-bound}
If $m\in12\ZZ+A$ is a nearest zero of $P$ to a real number $s$, then
\[
  P(s)\ge0.68(s-m)^2.
\]
\end{lemma}

\begin{proof}
On either half-gap from $m$, consider $P(m+u)/u^2$, with value $Q_m$ at
$u=0$.  Put $a=\pi/12$.  For a sine factor whose residue is $a'\ne m$
modulo $12$, the second derivative of its logarithm is
$-2a^2\csc^2(a(m+u-a'))<0$.  These factors stay nonzero on the half-gap.
The factor vanishing at $m$, after division by $u^2$, has the same
logarithmic concavity because
\[
 \frac{\dd^2}{\dd u^2}\log\abs{\frac{\sin(au)}{u}}
 =\frac1{u^2}-\frac{a^2}{\sin^2(au)}\le0.
\]
Here $\abs{\sin(au)}\le\abs{au}$, and the logarithm has a smooth limiting
value at $u=0$.  Squaring the divided sine only multiplies this second
derivative by two.
The quotient is therefore logarithmically concave on each half-gap,
including its continuous value at the center.  Its minimum is bounded
below by the smaller of its endpoint values.  At the center
$Q_m>1.75$.  At the gap midpoints, direct substitution gives the values
listed below; periodicity reduces the calculation to one period.
To evaluate them, put $\theta=\pi(s-2)/12$ and $c=\cos(2\theta)$.
The identity $4\sin v\sin w=2(\cos(v-w)-\cos(v+w))$ shows that the paired
unsquared factors for residues $(0,4)$, $(1,3)$, and $(7,9)$ are,
respectively,
\[
 1-2c,\qquad \sqrt3-2c,\qquad \sqrt3+2c.
\]
Consequently
\[
 P(s)=(4c^2-3)^2(1-2c)^2,
 \qquad c=\cos\frac{\pi(s-2)}6,
\]
which evaluates every entry directly.
\begin{center}
\begin{tabular}{c|c}
Gap between consecutive zeros & Midpoint quotient $P(s)/(s-m)^2$\\
\hline
$[0,1],\ [3,4]$ & $12-8\sqrt2$\\
$[1,3]$ & $1$\\
$[7,9]$ & $9$\\
$[4,7],\ [9,12]$ & $\dfrac89(3+\sqrt2+\sqrt3+\sqrt6)$
\end{tabular}
\end{center}
The distances to the two ends of a gap agree at its midpoint, so this
table supplies the endpoint value for both half-gaps.  The smallest value
is $12-8\sqrt2>0.68$: indeed $(283/200)^2>2$, so
$\sqrt2<283/200$ and $12-8(283/200)=0.68$.  The other three values are
at least $1$.  This proves the claim.
\end{proof}

For any $s\ge41.5$, choose a nearest node $m\ge43$; at $s=41.5$, choose
$43$ rather than the tied node $40$.  Both $H_i$ and $P R_{i,0}$ vanish
to order at least two at $m$, by \eqref{fourier:interpolation} and the absence of a
pole of $R_{i,0}$ there.  Thus
$E_i=H_i-P R_{i,0}$ vanishes to order at least two at $m$.  The segment between
$s$ and $m$ lies in $[41.5,\infty)$, and
\eqref{signs:tail-second-derivative} gives
\[
  \abs{E_i(s)}\le\frac{0.006}{2}(s-m)^2.
\]
Using \eqref{signs:rational-tail} and
Lemma~\ref{signs:P-lower-bound}, we obtain on the entire unbounded region
\begin{equation}\label{signs:tail-margin}
 \sigma_iH_i(s)
 \ge\left(0.011\cdot0.68-\frac{0.006}{2}\right)(s-m)^2
 =0.00448(s-m)^2\ge0.
\end{equation}

\begin{proposition}[Global signs and zeros]\label{signs:global}
The exact functions have the zero sets
\begin{equation}\label{signs:zero-sets}
 \{s\ge1:H_1(s)=0\}=\mathcal N,\qquad
 \{s\ge0:H_2(s)=0\}=\mathcal N,
\end{equation}
and satisfy the strict inequalities
\[
 H_1(s)<0\quad(s\ge1,\ s\notin\mathcal N),\qquad
 H_2(s)>0\quad(s\ge0,\ s\notin\mathcal N).
\]
The zero of $H_1$ at $1$ is simple.  Its zeros at $n\in\mathcal N$ with
$n>1$, and the zeros of $H_2$ at every $n\in\mathcal N$, have order exactly
two as zeros of the entire functions of $s$.
Consequently $F(x)\le0$ when $b\abs{x}^2\ge1$, and
$\widehat F(\xi)\ge0$ for every $\xi\in\RR^2$.
\end{proposition}

\begin{proof}
At every node $n\in\mathcal N$, the exact interpolation conditions give
$H_1(n)=H_2(n)=0$.  Away from the nodes, the quotient bounds
\eqref{signs:finite-quotient} are strict on the finite region, and
\eqref{signs:tail-margin} is strict on the remaining region.  This proves
\eqref{signs:zero-sets} and the asserted strict signs.

The interpolation slope $H_1'(1)=-Q_1\ne0$ makes its zero at $1$ simple.
For $H_2$ at every node in $\mathcal N\cap[1,40]$, and for $H_1$ at those
nodes greater than $1$, the chosen order is $\nu=2$.
The endpoint value in \eqref{signs:finite-quotient} is then
$\sigma_iH_i''(m)/2>0$, so the zero has order exactly two.
For a node $m\in\mathcal N$ with $m\ge43$, take $u$ sufficiently small that
$m$ remains the nearest node to $m+u$ and $m+u\ge41.5$.  The exact value and first derivative
at $m$ vanish, so \eqref{signs:tail-margin} yields
\[
 \frac{\sigma_iH_i''(m)}2
 =\lim_{u\to0}\frac{\sigma_iH_i(m+u)}{u^2}
 \ge0.00448>0.
\]
These zeros also have order exactly two.  Finally, the factors
$e^{-\pi h s}$ in $F$ and $\widehat F$ are positive, which gives their
non-strict sign inequalities.
\end{proof}
